% Still need to define the math: e.g. Let $\mathbb{G}$ be a cyclic group of prime order $q$ with generator $G$.
\section{Proposed Scheme}
\label{sec:scheme}

We propose a credential mechanism for destination-based authorization in mixnets, integrated into a variant of the Sphinx packet format. 
The objective is to allow each mixnode to verify that the destination carried in the encrypted routing information has been authorized, without the verification itself disclosing that destination. 

Figure~\ref{fig:architecture} summarizes the interactions between the client, the credential authorities, and the mixnet. 
Following system initialization, the protocol comprises four stages:
\begin{enumerate}
    \item \emph{Credential request.}
    The client submits a blinded destination point to a set of authorities,
    together with a zero-knowledge proof that it encodes an address outside
    a publicly available blacklist.
    \item \emph{Partial credential issuance.}
    Each contacted authority verifies the proof and, if verification succeeds,
    returns a \emph{credential share}.
    \item \emph{Credential aggregation.}
    The client combines a threshold number of shares and unblind the
    destination to obtain a \emph{destination-bound credential}.
    \item \emph{Header binding.}
    During packet construction, the client rerandomizes the destination-bound
    credential and bound it with the routing information to produce a
    \emph{header-bound credential}.
\end{enumerate}

\begin{figure}[t]
    \centering
    \resizebox{\linewidth}{!}{\newcommand{\WIDTH}{10 mm}
\newcommand{\HEIGHT}{18 mm}

\tikzset{
  mixnode/.style={draw, circle, radius=2mm},
  authority/.style={draw, inner sep=0pt, minimum width=1cm, minimum height=1cm},
}

\newcommand{\client}[1]{
  \begin{tikzpicture}[scale=0.4]
    % Monitor screen
    \draw[fill=#1!20, draw=black, rounded corners=4pt] (0,0) rectangle (4,2.5);
    % Inner bezel
    \draw[fill=white, draw=none, rounded corners=3pt] (0.2,0.2) rectangle (3.8,2.3);
    % Stand
    \draw[fill=#1!40, draw=black] (1.75,-0.7) rectangle (2.25,0);
    % Base
    \draw[fill=#1!40, draw=black, rounded corners=2pt] (0.8,-1) rectangle (3.2,-0.7);
  \end{tikzpicture}
}


\begin{tikzpicture}

% Nodes
\node (client) [label=below:\text{\color{gray}Client}] at (0, 0) {\scalebox{1}{\client{gray}}};

\node (Authority1) [authority, color=black] at ($(client.west) + (-2.5*\WIDTH, 0.8*\HEIGHT)$)  {};
\node (Authority2) [authority, color=black] at ($(client.west) + (-2.5*\WIDTH, 0.2*\HEIGHT)$) {};
\node (Authority3) [authority, color=black] at ($(client.west) + (-2.5*\WIDTH, -0.8*\HEIGHT)$)  {};
\draw[thick, dashed, outer sep=5mm, shorten <= 2mm, shorten >= 2mm] (Authority2) -- (Authority3);
\node (AuthorityBox) [draw, densely dashed, color=gray, fit={(Authority1) (Authority3)}, rounded corners=5pt, inner sep=6pt, label=above:\text{\color{gray}Credential Authorities}] {};

\node (entry) [draw, minimum height=2cm, minimum width=0.4cm] at ($(client.east) + (2*\WIDTH, 0)$) {};
\node[align=center, font=\scriptsize, rotate=90] at (entry.center) {Entry gateway};
\node (11) [mixnode] at ($(entry.north east) + (\WIDTH/1.5, -\HEIGHT)$) {};
\node (12) [mixnode] at ($(entry.east) + (\WIDTH/1.5, 0)$) {};
\node (13) [mixnode] at ($(entry.south east) + (\WIDTH/1.5, \HEIGHT)$) {};
\node (21) [mixnode] at ($(11.east) + (\WIDTH/1.5, 0)$) {};
\node (22) [mixnode] at ($(12.east) + (\WIDTH/1.5, 0)$) {};
\node (23) [mixnode] at ($(13.east) + (\WIDTH/1.5, 0)$) {};
\node (31) [mixnode] at ($(21.east) + (\WIDTH/1.5, 0)$) {};
\node (32) [mixnode] at ($(22.east) + (\WIDTH/1.5, 0)$) {};
\node (33) [mixnode] at ($(23.east) + (\WIDTH/1.5, 0)$) {};
\node (exit) [draw, minimum height=2cm, minimum width=0.4cm] at ($(32.east) + (\WIDTH/1.5, 0)$) {};
\node[align=center, font=\scriptsize, rotate=90] at (exit.center) {Exit gateway};
\node (mixnet) [draw, densely dashed, color=gray, fit={(entry) (exit)}, rounded corners=5pt, inner sep=10pt, label=above:\text{\color{gray}Mix Network}, label=below:\text{\color{gray}Credential verification (at any mixnode)}] {};


\draw[-Latex, shorten >= 1mm] ($(client.west) + (0, 4mm)$) to node[font=\small, midway, fill=white, inner sep=1pt] {\color{black} (1)} ($(AuthorityBox.east) + (0, 4mm)$);
\draw[-Latex, shorten <= 1mm] ($(AuthorityBox.east) + (0, -4mm)$) to node[font=\small, midway, fill=white, inner sep=1pt] {\color{black} (2)} ($(client.west) + (0, -4mm)$);
\draw[-Latex] (client) to[out=120,in=60,looseness=2] node[font=\small, midway, fill=white, inner sep=1pt] {\color{black} (3)} (client);
\draw[-Latex, shorten >= 1mm] (client.east) to node[font=\small, midway, fill=white, inner sep=1pt] {\color{black} (4)} (mixnet.west);

\draw (entry.east) -- (11.west);
\draw (entry.east) -- (12.west);
\draw (entry.east) -- (13.west);
\draw (11.east) -- (21.west);
\draw (11.east) -- (22.west);
\draw (11.east) -- (23.west);
\draw (12.east) -- (21.west);
\draw (12.east) -- (22.west);
\draw (12.east) -- (23.west);
\draw (13.east) -- (21.west);
\draw (13.east) -- (22.west);
\draw (13.east) -- (23.west);
\draw (21.east) -- (31.west);
\draw (21.east) -- (32.west);
\draw (21.east) -- (33.west);
\draw (22.east) -- (31.west);
\draw (22.east) -- (32.west);
\draw (22.east) -- (33.west);
\draw (23.east) -- (31.west);
\draw (23.east) -- (32.west);
\draw (23.east) -- (33.west);
\draw (31.east) -- (exit.west);
\draw (32.east) -- (exit.west);
\draw (33.east) -- (exit.west);

\end{tikzpicture}}
    \caption{Overview of the credential issuance and header binding. 
    (1) After proving that a \emph{blinded destination} satisfies the authorization policy, 
    (2) the client obtains \emph{credential shares} from a threshold set of authorities. 
    (3) It aggregates and unblinds these shares to obtain a \emph{destination-bound credential}.
    (4) When sending a packet, a \emph{header-bound credential} is produced by rerandomizing the credential and bound it to the encrypted routing information.}
    \label{fig:architecture}
\end{figure}


\paragraph{Notation}
Elliptic curve points are denoted by large Greek letters, while calligraphic letters denote ordered collections of points, such as $\mathcal{H}_i$, $\mathcal{L}_i$, and $\mathcal{B}_i$.
Let $\mathbb{G}_1$ and $\mathbb{G}_2$ be additive cyclic groups of prime order $q$, with respective generators $\ZETA_1$ and $\ZETA_2$, 
and let $e:\mathbb{G}_1\times\mathbb{G}_2\rightarrow\mathbb{G}_T$ be a non-degenerate type-3 bilinear pairing. 
Scalars belong to $\mathbb{F}_q$, and $[s]\RHO$ denotes the scalar multiplication of a point $\RHO$ by $s$. 
Hop indices are denoted by $i$, vector indices by $j$, and authority indices by $k$.
A tilde denotes a blinding point such as $\blind{\RHO} = [r]\RHO$, where $r$ is a random scalar.
Column vectors represent collection of element of the same type, while tuple components have different types.


\paragraph{DSphinx header}
The proposed scheme builds on DSphinx~\cite{DSphinx}, an adaptation of Sphinx in which each header element is represented as an elliptic-curve point.
At hop $i$, the header $\mathcal{H}_i$ has the structure illustrated in Fig.~\ref{fig:header}.
The point $\ALPHA_i$ allows mixnode $i$ to derive its shared secret using its secret key.
For a route of $\ell$ mixnodes, the routing layer $\mathcal{L}_i$ is considered as a column vector of $2\ell+1$ points in $\mathbb{G}_1$, 
where $\ETA_i$ is the next-hop identifier , $\GAMMA_i$ is the integrity tag, and $\mathcal{B}_i$ is the encrypted routing information represented as a vector of $2\ell-1$ points.
$\SIGMA_i$ is the \emph{header-bound} credential.
All point components belong to $\mathbb{G}_1$.
\begin{equation}\label{eq:header}
\end{equation}
\begin{figure}
    \centering
    $$
    \mathcal{H}_i=(\ALPHA_i, \mathcal{L}_i, \SIGMA_i),
    \qquad
    \mathbf{L}_i=
    \begin{bmatrix}
        \ETA_i \\ \GAMMA_i\\ \mathcal{B}_i,
    \end{bmatrix}
    \qquad
    \mathcal{B}_i=
    \begin{bmatrix}
        \BETA_1 \\ \vdots \\ \BETA_{2\ell-1},
    \end{bmatrix}.
    $$
    \resizebox{1.\linewidth}{!}{\centering
\begin{tikzpicture}

    \node (H) at (0,2) {$\mathcal{H}_i = $};
    \node (alpha) [block=1] at (H.east) {$\ALPHA_{i}$};
    \node (N) [block=1] at (alpha.east) {$\ETA_i$};
    \node (gamma) [block=1] at (N.east) {$\GAMMA_{i}$};
    \node (Bi) [block=5] at (gamma.east) {$\mathcal{B}_i$};
    \node (Sigma) [block=1] at (Bi.east) {$\SIGMA_{i}$};

    \node (B) [block=5, anchor=south west, yshift=5pt] at (Bi.north west)  {};
    \foreach \j in {1,...,4} {
        \draw[dotline] ($(B.north west) + (\j*\width, 0)$) -- ($(B.south west) + (\j*\width, 0)$) {}; 
    }
    \node (beta_1) [blockspace=1] at (B.west)  {$\beta_{\scriptscriptstyle 1}$};
    \node (beta_2) [blockspace=1] at (beta_1.east)  {$\ldots$};
    \node (beta_3) [blockspace=1] at (beta_2.east)  {$\beta_{\scriptscriptstyle j}$};
    \node (beta_4) [blockspace=1] at (beta_3.east)  {$\ldots$};
    \node (beta_5) [blockspace=1] at (beta_4.east)  {$\beta_{\scriptscriptstyle 2\ell-1}$};
    \node[anchor=east] at (B.west) {$\mathcal{B}_i = $};

    % Notes
    \draw [decorate, decoration={brace, amplitude=3pt, mirror, raise=4pt}] (N.215) -- (Bi.south east) node[midway, below=8pt] {$\mathcal{L}_i$};

\end{tikzpicture}}
    \caption{Header structure at hop $i$.}
    \label{fig:header}
\end{figure}

\paragraph{Address representations}
A public directory associates each mixnode identifier $\ETA_i$ with its network address $\mathrm{IP}_i$, allowing address recovery by lookup.
Mixnode address $\mathrm{IP}_i$ is encoded by multiplying $\ZETA_1$ by a publicly computable scalar:
$$ \ETA_i=[h_{\mathbb{F}_q}(\mathrm{IP}_i)]\ZETA_1, \qquad h_{\mathbb{F}_q}:\mathbb{Z}\to\mathbb{F}_q. $$

For destinations, maintaining such a directory would be impractical, and the encoding must not reveal the scalar relating the destination point to $\ZETA_1$.
We therefore use a reversible \emph{try-and-increment} coordinate embedding.
For a destination address $d$, we consider $x_d=d\,2^\kappa+\tau+c$, where $\tau$ is a publicly prescribed $\kappa$-bit padding value (e.g. zeroes or a truncated timestamp).
Starting from zero, the counter $c$ is incremented until a point $(x_d,y_d)\in\mathbb{G}_1$ is found, subject to $\tau+c<2^\kappa$ and $x_d<p$, where $\mathbb{F}_p$ is the curve's coordinate field.
The coordinate $y_d$ is selected using a fixed sign convention.
The padding reserves space for the increment while preserving the address bits. 
The padding and point-selection rules are fixed by the protocol and enforced by the zero-knowledge circuit.
The encoding and decoding functions are thus defined by:
$$ \operatorname{Enc}(d)=(x_d,y_d)=\D{}, \qquad \operatorname{Dec}(\D{}) =\left\lfloor\frac{x_d}{2^\kappa}\right\rfloor. $$ %  \qquad \operatorname{Dec}(\operatorname{Enc}(d)) = d
\todo[inline]{Counter is bounded}

\subsection{Threshold Multi-Authority Setup}
\label{sec:scheme-setup}

Let $\mathcal{A}=\{A_1,\ldots,A_{n_A}\}$ be a set of $n_A$ credential authorities, with issuance threshold $2\leq t\leq n_A$. 
Each authority $A_k$ has a public identifier $x_k\in\mathbb{F}_q$.

At initialization, the authorities are assumed to execute a secure distributed key-generation (DKG) protocol~\cite{DKG}. 
It outputs a Shamir sharing of a secret signing scalar $y$~\cite{ShamirSS}. 
In particular, there exists a polynomial $f\in\mathbb{F}_q[z]$ of degree at most $t-1$ such that:
$$ y=f(0), \qquad y_k=f(x_k),$$
where authority $A_k$ holds only its secret share $y_k$.
The authorities publish the public key $\PSI_1=[y]\ZETA_1$ and the system verification key $\PSI_2=[y]\ZETA_2$ using Lagrange interpolation 
of their shares $[y_k]\ZETA_1$ and $[y_k]\ZETA_2$ from any set $\mathcal{I}$ of at least $t$ authorities, without reconstructing $y$.
$$
\PSI_i = \sum_{k\in\mathcal{I}} [\lambda_k^{\mathcal{I}}]\bigl([y_k]\ZETA_i\bigr) = [y]\ZETA_i, 
\qquad \lambda_k^{\mathcal{I}} = \prod_{\substack{j\in\mathcal{I}\\j\neq k}} \frac{-x_j}{x_k-x_j}.
$$

The system also establishes public parameters $\mathsf{pp}_{\mathrm{zk}}$ for a non-interactive argument of knowledge. 
We instantiate the proof system with Groth16, using a Powers of Tau ceremony followed by a circuit-specific setup to generate the common reference string (CRS) and the associated proving and verification keys. 
Since the authorization circuit is fixed and the blacklist is public, these parameters are generated once and reused across credential requests.


\subsection{Credential Issuance} 
\label{sec:scheme-issuance}
During credential issuance, the authorities require a proof that the requested destination $d$ is not included in the public blacklist $\mathbb{B}$, without learning $d$.
Credential lifetime and subsequent policy changes require a separate validity or revocation mechanism that is outside the scope of this work.

\paragraph{Zero-knowledge Proof}
\label{sec:zk-proof}
The client encode the desired destination $d$ into an elliptic curve point $\D{}=\operatorname{Enc}(d)$, 
samples a fresh blinding scalar $b \in \mathbb{F}_q$, and computes $ \widetilde{\D{}}=[b]\D{}. $
It then generates a NIZK proof $\pi$ for the constraints:
\begin{equation}
    % \label{eq:authorization-relation}
        \begin{aligned}
        &\pi \bigl((\widetilde{\D{}},\mathbb{B});(d,b)\bigr)=1 \quad\Longleftrightarrow\quad
        \begin{cases}
            d\in\{0,\ldots,2^{128}-1\},\\
            d\notin\mathbb{B},\\
            \widetilde{\D{}}
            =[b]\operatorname{Enc}(d).
        \end{cases}
    \end{aligned}
\end{equation}
The public statement consists of the blinded destination $\widetilde{\D{}}$ and the blacklist $\mathbb{B}$.
The destination address $d$ and blinding scalar $b$ form the private witness.
The first condition verify that private destination is in the range of an IP address.
The second condition verify the non-membership to the blacklist.
The last condition binds the non-membership proof to the blinded destination submitted for credential issuance.

\paragraph{Partial Credential Issuance}
\label{sec:scheme-partial-issuance}
The client sends $(\widetilde{\D{}},\pi)$ to at least $t$ authorities.
Each authority verifies the proof $\pi$ using $\mathsf{pp}_{\mathrm{zk}}$ and the provided blinded encoded destination.
A request is rejected if the proof verification fails.
After accepting the request, authority $A_k$ computes $\SIGMA_{k} = [y_k] \widetilde{\D{}}$, and returns the credential share $\SIGMA_{k}$ to the client.


\paragraph{Credential Aggregation}
\label{sec:scheme-aggregation}

After collecting shares from a set $\mathcal{I}$ of at least $t$ distinct authorities, the client aggregates the credential shares using Lagrange interpolation.
It then unblinds the credential by multiplying it by $b^{-1}$. 
This yields the \emph{destination-bound} credential $\SIGMA_{\Delta} = [y] \Delta$.
An honestly generated credential satisfies $e(\SIGMA_{\Delta},\ZETA_2)=e(\Delta,\PSI_2)$.
During issuance, destination blinding and the zero-knowledge proof $\pi$ allow the authorities to verify that the requested destination satisfies the authorization policy without learning the desired destination.




\subsection{Header Construction}
\label{sec:header-construction}

The client selects a route, derives the shared secrets, and constructs the routing layers from the last hop to the first as shown in Figure~\ref{fig:header-construction}.
\begin{figure}
    \centering
    \resizebox{0.6\linewidth}{!}{\centering
\begin{tikzpicture}
    % \tiny
    \setlength{\y}{0cm}
    \foreach \j in {0,...,4} {
        \node (b3\j) [blockspace=1] at (\j*\width, \y) {};
        \ifnum\j>1  \draw[dotline, gray] (b3\j.north west) -- (b3\j.south west);\fi
    }
    \node (dest) [block=1] at (0, \y) {$s_3^{-1}\Delta$};
    \node (fillers) [block=4] at (\width, \y) {$\PHI$};  
    \node[left of=dest] {$\beta_3 =$};

    % HMAC 3
    \vGap
    \node (hmac3) [eq] at (1.5*\width, \y) {\ref{eq:integrity}};
    \node[left=\width of hmac3] (input_hmac3) {$s_{3}$};
    \vGap
    \node (L3) [blockspace=7] at (0, \y) {};
    \foreach \j in {0,...,6} {\node (L3\j) [blockspace=1] at (\j*\width, \y) {}; }
    \foreach \j in {0,...,4} {
        \node (beta3\j) [blockspace=1] at (2*\width + \j*\width, \y) {}; 
        \ifnum\j>0  \draw[dotline, gray] (beta3\j.north west) -- (beta3\j.south west);\fi
    }
    \node (N3) [block=1] at (0, \y) {$\ETA_{3}$}; 
    \node (gamma3) [block=1] at (\width, \y) {$\GAMMA_3$}; 
    \node (beta3) [block=5] at (2*\width, \y) {$\beta_3$};  
    \node[left of=N3] {$\mathcal{L}_3 =$};

    % ENC 2
    \vGap
    \node [ovale=5] at (0, \y) {$s_2$};
    \vGap
    \node [block=5] at (0, \y) {$\beta_2$};
    \foreach \j in {0,...,4} {
        \node (b2\j) [blockspace=1] at (\j*\width, \y) {}; 
        \ifnum\j>0  \draw[dotline, gray] (b2\j.north west) -- (b2\j.south west);\fi
    }

    % HMAC2
    \vGap
    \node (hmac2) [eq] at (1.5*\width, \y) {\ref{eq:integrity}};
    \node[left=\width of hmac2] (input_hmac2) {$s_{2}$};
    \vGap
    \node (L2) [blockspace=7] at (0, \y) {};
    \foreach \j in {0,...,6} {\node (L2\j) [blockspace=1] at (\j*\width, \y) {}; }
    \foreach \j in {0,...,4} {
        \node (beta2\j) [blockspace=1] at (2*\width + \j*\width, \y) {}; 
        \ifnum\j>0  \draw[dotline, gray] (beta2\j.north west) -- (beta2\j.south west);\fi
    }
    \node (N2) [block=1] at (0, \y) {$\ETA_{2}$}; 
    \node (gamma2) [block=1] at (\width, \y) {$\GAMMA_2$}; 
    \node (beta2) [block=5] at (2*\width, \y) {$\beta_2$};  
    \node[left of=N2] {$\mathcal{L}_2 =$};

    % ENC 1
    \vGap
    \node [ovale=5] at (0, \y) {$s_1$};
    \vGap
    \node [block=5] at (0, \y) {$\beta_1$};
    \foreach \j in {0,...,4} {
        \node (b1\j) [blockspace=1] at (\j*\width, \y) {}; 
        \ifnum\j>0  \draw[dotline, gray] (b1\j.north west) -- (b1\j.south west);\fi
    }

    % HMAC1
    \vGap
    \node (hmac1) [eq] at (1.5*\width, \y) {\ref{eq:integrity}};
    \node[left=\width of hmac1] (input_hmac1) {$s_{1}$};
    \vGap
    \node (L1) [blockspace=7] at (0, \y) {};
    \foreach \j in {0,...,6} {\node (L1\j) [blockspace=1] at (\j*\width, \y) {}; }
    \foreach \j in {0,...,4} {
        \node (beta1\j) [blockspace=1] at (2*\width + \j*\width, \y) {}; 
        \ifnum\j>0  \draw[dotline, gray] (beta1\j.north west) -- (beta1\j.south west);\fi
    }
    \node (N1) [block=1] at (0, \y) {$\ETA_{1}$}; 
    \node (gamma1) [block=1] at (\width, \y) {$\GAMMA_1$}; 
    \node (beta1) [block=5] at (2*\width, \y) {$\beta_1$};  
    \node[left of=N1] {$\mathcal{L}_1 =$};

    %% ARROWS %%
    % hmac 3
    \begin{pgfonlayer}{background}
        \foreach \j in {0,...,4} {
            \draw[->, dashed, very thin, color=unHighlight!80] (b3\j.south) -- ([yshift=+1pt] beta3\j.north);
        }
    \end{pgfonlayer}
    \draw[arrow] (input_hmac3) -- (hmac3.west);
    \draw[arrow] (hmac3.south) -- (gamma3.north);
    \draw [decorate, decoration={brace, amplitude=3pt, mirror, aspect=0.3, raise=2pt}] (dest.south west) -- (fillers.south east);
    % ENC 2
    \begin{pgfonlayer}{background}
        \foreach \j in {0,...,4} {
            \draw[->, dashed, very thin, color=unHighlight!80] (L3\j.south) -- ([yshift=+1pt] b2\j.north);
        }
    \end{pgfonlayer}
    % hmac 2
    \begin{pgfonlayer}{background}
        \foreach \j in {0,...,4} {
            \draw[->, dashed, very thin, color=unHighlight!80] (b2\j.south) -- ([yshift=+1pt] beta2\j.north);
        }
    \end{pgfonlayer}
    \draw[arrow] (input_hmac2) -- (hmac2.west);
    \draw[arrow] (hmac2.south) -- (gamma2.north);
    \draw [decorate, decoration={brace, amplitude=3pt, mirror, aspect=0.3, raise=2pt}] (b20.south west) -- (b24.south east);
    % ENC 1
    \begin{pgfonlayer}{background}
        \foreach \j in {0,...,4} {
            \draw[->, dashed, very thin, color=unHighlight!80] (L2\j.south) -- ([yshift=+1pt] b1\j.north);
        }
    \end{pgfonlayer}
    % hmac 1
    \begin{pgfonlayer}{background}
        \foreach \j in {0,...,4} {
            \draw[->, dashed, very thin, color=unHighlight!80] (b1\j.south) -- ([yshift=+1pt] beta1\j.north);
        }
    \end{pgfonlayer}
    \draw[arrow] (input_hmac1) -- (hmac1.west);
    \draw[arrow] (hmac1.south) -- (gamma1.north);
    \draw [decorate, decoration={brace, amplitude=3pt, mirror, aspect=0.3, raise=2pt}] (b10.south west) -- (b14.south east);

   
\end{tikzpicture}
}
    \caption{DSphinx header construction for a path of length $\ell=3$ (without credential binding). 
    Circled numbers \eq[\color{red}x] refers to Equation x in the paper.}
    \label{fig:header-construction}
\end{figure}


\paragraph{Route and shared secrets}
The client samples an ordered route $(\ETA_1,\ldots,\ETA_\ell)$, a nonce $x_1\in\mathbb{F}_q$, and compute the cryptographic element $\ALPHA_1 = [x_1]\ZETA_1$.
For each hop $i=1,\ldots,\ell$, it uses the mixnode public key $\KAPPA_i=[\mathsf{k}_i]\ZETA_1$ to derive the corresponding shared secrets in cascade as:
\begin{equation}
    % \label{eq:shared_secret}
    \{s_i\}_{i=1}^{\ell} \quad\left\{\quad
    \begin{aligned}
        s_i &= h_{\mathbb{F}_q}\bigl([x_i]\KAPPA_i\bigr),\\
        x_{i+1} &= s_i x_i.
    \end{aligned}
    \right.
\end{equation}

\paragraph{Preprocessing}
The client prepares filler vector $\PHI$ that maintain the routing-layer size during processing. 
Starting with the empty vector $\PHI^{(0)}$, it applies the following recurrence for $i=1,\ldots,\ell-1$, obtaining a filler vector $\PHI^{\ell-1}$ with $2(\ell-1)$ elements in $\mathbb{G}_1$. 
\begin{equation}
    % \label{eq:preprocessing}
    \PHI^{(i)} = [s_i]
    \begin{bmatrix}
        \PHI^{(i-1)}\\
        h_{\mathbb{G}}(s_i)\\
        \ETA_i
    \end{bmatrix}
    \qquad (i=1,\ldots,\ell-1).
\end{equation}

\paragraph{Layer construction.}
The client constructs the routing layers in reverse order, from the last hop $i=\ell$ to the first hop $i=1$.
It initializes the innermost routing vector $\mathcal{B}_\ell$ by concatenating the destination point multiplied by the inverse of the last hop shared secret $[s_\ell^{-1}]\D{}$, with the precomputed filler vector  $\PHI^{\ell-1}$.
\begin{equation}
    % \label{eq:inner_layer}
    \mathcal{B}_\ell = 
    \begin{bmatrix}
        [s_\ell^{-1}]\D{}\\
        \phi^{(\ell-1)}
    \end{bmatrix}
\end{equation}
For each preceding hop $i=\ell-1,\ldots,1$, it forms $\mathcal{B}_i$ by discarding the last two elements of $\mathcal{L}_{i+1}$ and multiplying the resulting vector by $s_i^{-1}$.
For every layer, the client computes an integrity tag $\GAMMA_i$ over $\mathcal{B}_i$ using the shared secret $s_i$, then prepends the hop identifier $\ETA_i$ and the tag to obtain $\mathcal{L}_i$.
\begin{equation}
    \label{eq:layer}
    \mathcal{B}_i = [s_i^{-1}] (\mathcal{L}_{i+1})_{1:2\ell-1},
    \qquad
    \GAMMA_i = \mathsf{HMAC}(s,\mathcal{B}_i),
    \qquad
    \mathcal{L}_i = \begin{bmatrix} \ETA_i \\ \GAMMA_i \\ \mathcal{B}_i \end{bmatrix}.
\end{equation}


\subsection{Credential Binding}
\label{sec:credential-binding}
The client can link the \emph{destination-bound} credential $\SIGMA_{\D{}}$ to the encrypted routing information from the header to obtain a \emph{header-bound} credential $\SIGMA_1$, 
as in eq.~\ref{eq:header_bound_credential}, where $w_i=\prod_{j=1}^{i-1}s_j$, with $w_1=1$, and $IP_i$ is the address of hop $i$.
The client can then send the packet with the following header $\mathcal{H}_1=(\ALPHA_{1},\mathcal{L}_1,\SIGMA_1)$ to the first mixnode.
\begin{equation}
\label{eq:header_bound_credential}
    \SIGMA_{1} = [w_{\ell+1}^{-1}] \SIGMA_{\D{}} + \sum_{i=1}^{\ell}[w_i^{-1} h_{\mathbb{F}_q}\bigl(IP_i\bigr)] \PSI_1,\\
\end{equation}

\subsection{Packet Processing \& Credential Verification}
\label{sec:packet-processing}

Upon receiving header $\mathcal{H}_i=(\ALPHA_i,\mathcal{L}_i,\SIGMA_i)$, mixnode $i$ verifies the credential and the integrity tag before processing and forwarding the header as shown in Figure~\ref{fig:header-processing}.
\begin{figure}
    \centering
    \resizebox{1.\linewidth}{!}{\centering
\begin{tikzpicture}
    %% INPUT
    \node (L0) at (0,0) {};
    \node (A) [block=1] at (L0) {$\ALPHA_{i}$};
    \node (N) [block=1] at ($(A.east) + (0.5*\width, 0)$) {$\ETA_i$};
    \node (C) [block=1] at ($(N.east) + (0.1*\width, 0)$) {$\GAMMA_{i}$};
    \node (B) [block=5] at ($(C.east) + (0.5*\width, 0)$)  {$\mathcal{B}_{i}$};
    \node (Sign) [block=1] at ($(B.east) + (2*\width, 0)$) {$\SIGMA_{i}$};

    \foreach \j in {1,...,4} {
        \draw[dotline] ($(B.north west) + (\j*\width, 0)$) -- ($(B.south west) + (\j*\width, 0)$) {}; 
    }

    %% LINE 1 (shared secret, check integrity & decrypt)
    \node (L1) at ($(L0) - (0, 2.5*\height)$) {};
    \node (sk) [block=1, color=red] at ($(L1) - (1.25*\width, 0)$) {\sk{i}};
    \node (ss) [eq] at (A |- L1) {\ref{eq:derive_secret}};
    \node (S) [block=1] at (N.west |- L1)  {\s{i}};
    \node (check) [eq] at (C |- L1)  {\ref{eq:integrity_check}};
    \node (credEq) [eq] at (Sign.-110 |- L1) {\ref{eq:credential_verification}};
 
    %% LINE 2 (decryption up)
    \node (L2) at ($(L1) - (0, 1*\height)$) {};
    \node (BB) [block=5] at (B.west |- L2) {$\mathcal{B}_i$};
    \foreach \j in {1,...,4} {
        \draw[dotline] ($(BB.north west) + (\j*\width, 0)$) -- ($(BB.south west) + (\j*\width, 0)$) {}; 
    }
    \node (pad) [zero_pad=1] at (BB.east) {$h(s_i)$}; 
    \node () [zero_pad=1] at (pad.east) {$\ETA_i$}; 
    \node (xor) [eq] at ($(check |- BB) - (0, 1*\height)$) {\ref{eq:decrypt}};

    %% LINE 3 (decryption down)
    \node (L3) at ($(L2) - (0, 2*\height)$) {};
    \node (N'_) [block=1] at (BB.west |- L3) {$\ETA_{i+1}$};
    \node (C'_) [block=1] at (N'_.east) {$\GAMMA_{i+1}$};
    \node (B'_) [block=5] at (C'_.east) {$\mathcal{B}_{i+1}$};
    \foreach \j in {1,...,4} {
        \draw[dotline] ($(B'_.north west) + (\j*\width, 0)$) -- ($(B'_.south west) + (\j*\width, 0)$) {}; 
    }

    %% LINE 4 (update)
    \node (L4) at ($(L2) - (0, 3*\height)$) {};
    \node (Si) [block=0.5, anchor=center, lightgray] at (credEq |- xor) {$\XI_i$};
    \node (blind) [eq] at (A.-60 |- L4) {\ref{eq:update_alpha}};
    \node (updSign) [eq] at (credEq |- L4) {\ref{eq:update_credential}};
    
    %% BOX
    \node (mixnode) [yshift=1pt, inner ysep=7pt, thick, draw=red, dashed, fit=(sk) (updSign.south-|Sign.east)] {};
    \node[xshift=-7pt, rotate=90] at (mixnode.west) {\color{red} \textbf{Mixnode $i$}};

    %% OUTPUT
    \node (L) at ($(mixnode.south) - (0, 1.5*\height)$) {};
    \node (A') [block=1] at (A.west |- L) {$\ALPHA_{i+1}$};
    \node (N') [block=1] at (N.west |- L) {$\ETA_{i+1}$};
    \node (C') [block=1] at (C.west |- L) {$\GAMMA_{i+1}$};
    \node (B') [block=5] at (B.west |- L)  {$\mathcal{B}_{i+1}$};
    \foreach \j in {1,...,4} {
        \draw[dotline] ($(B'.north west) + (\j*\width, 0)$) -- ($(B'.south west) + (\j*\width, 0)$) {}; 
    }
    \node (Sign') [block=1] at (Sign.west |- L) {$\SIGMA_{i+1}$};
    %% ARROW
    % (4) sum chunks
    \draw[arrow, lightgray] ($(BB.south) + (-2*\width, 0)$) |- (Si);
    \draw[arrow, lightgray] ($(BB.south) + (0*\width, 0)$) |- (Si);
    \draw[arrow, lightgray] ($(BB.south) + (2*\width, 0)$) |- (Si);
    \draw[arrow, lightgray] ($(BB.south) + (4*\width, 0)$) |- (Si);
    % (5) credential verification
    \draw[arrow] (Sign.-40) |- (credEq);
    \draw[arrow, lightgray] (Si) -- (credEq);
    (credEq);
    % (6) shared secret
    \draw[arrow] (A.south -| ss) -- (ss);
    \draw[arrow] (sk) -- (ss);
    \draw[arrow] (ss) -- (S);
    % (7) check
    \draw[arrow] (B.south) |- (check);
    \draw[arrow] (C) -- (check);
    \draw[arrow] (S.east |- check) -- (check); 
    % (8) xor
    \draw[arrow] (B.south) -- (BB.north);
    \draw[arrow] (S.south) |- (xor);
    \draw[arrow] (BB.west) -| (xor);
    \draw[arrow] (xor) |- (N'_.west);
    % (9) blind
    \draw[line] (A.south -| blind) -- (ss -| blind);
    \draw[arrow] (ss -| blind) -- (blind);
    \draw[arrow] (S) |- (blind);
    \draw[arrow] (blind) -- (blind |- A'.north);
    % (10) update credential
    \draw[arrow] (S) |- (updSign);
    \draw[arrow] (Sign.-40) |- (updSign);
    % \draw[arrow, lightgray] (Si.south) -| (updSign);
    \draw[arrow] (updSign.south) -- ($(updSign.south) - (0, 0.9*\height)$) -| (Sign'.110);
    % output
    \draw[arrow] (N'_) -- ($(N'_.south) - (0, 1.7*\height)$) -| (N');
    \draw[arrow] (C'_) -- ($(C'_.south) - (0, 1.9*\height)$) -| (C');
    \draw[arrow] (B'_.south) -- ($(B'_.south) - (0, 1.8*\height)$) -| (B');
\end{tikzpicture}}
    \caption{DSphinx header and its credential processed by mixnode $i$. 
    Circled numbers \eq[\color{red}x] refers to Equation x in the paper.}
    \label{fig:header-processing}
\end{figure}

\paragraph{Credential verification}
The node computes $\XI_i$ by summing the odd-indexed elements of the layer vector $\mathcal{L}_i$, which contains all the route and destination information.
\begin{equation}\label{eq:sum_route}
    \XI_i 
    \;=\; \sum_{j=0}^{\ell}(\mathcal{L}_i)_{2j+1} 
    \;=\; \ETA_i + \sum_{j=1}^{\ell}\BETA_{2j-1}
\end{equation}
It then verifies the pairing equation~\eqref{eq:credential_verification} and discards the packet if the check fails.
Importantly, this verification does not require identifying the vector element containing the destination, preventing the mixnode from determining its position in the route (except at the first and last hops).
\begin{equation}
\label{eq:credential_verification}
    e(\SIGMA_{i}, \ZETA_2)
    \overset{?}{=}
    e(\XI_i, \PSI_2)%,
\end{equation}

\paragraph{Integrity check}
The mixnode derives the shared secret $s_i$ using its secret key $k_i$ with the cryptographic element $\ALPHA_i$ from the header.
Then verify the integrity tag using eq~\eqref{eq:integrity}. 
The packet is discarded if integrity check fails.
\begin{align}
    s_i &=   h_{\mathbb{F}_q}\bigl([k_i]\ALPHA_i \bigr), \label{eq:derive_secret}\\
    \GAMMA_i & \overset{?}{=} \mathsf{HMAC}(s,\mathcal{B}_i). \label{eq:integrity_check}
\end{align}

\paragraph{Layer processing}
After both checks succeed, the next layer $\mathcal{L}_{i+1}$ is processed as in eq.\ref{eq:decrypt}. 
The two firsts elements of $\mathcal{L}_{i+1}$ reveals the next routing identifier $\ETA_{i+1}$ and its integrity tag $\GAMMA_{i+1}$.
While the remaining forms the next encrypted routing information $\mathcal{B}_{i+1}$.
\begin{equation}
    \label{eq:decrypt}
    \mathcal{L}_{i+1} = [s_i] \begin{bmatrix} \mathcal{B}_i \\ h_{\mathbb{G}}(s_i) \\ \ETA_i \end{bmatrix}
\end{equation}

\paragraph{Header update}
The node updates the cryptographic element $\ALPHA$ and credential $\SIGMA$ as in \eqref{eq:update_alpha}  and \eqref{eq:update_credential}.
Since the layer decryption gives $\XI_{i+1}=[s_i]\XI_i$, these updates preserve the pairing equation: 
$ e(\sigma_{i+1},\zeta_2) = e([s_i]\sigma_i,\zeta_2) = e(\xi_i,\psi_2)^{s_i} = e([s_i]\xi_i,\psi_2) = e(\xi_{i+1},\psi_2). $
\begin{equation}
    \label{eq:update_alpha}
    \ALPHA_{i+1} = [s_i]\ALPHA_{i} 
\end{equation}
\begin{equation}
    \label{eq:update_credential}
    \SIGMA_{i+1} = [s_i] \SIGMA_{i} %\bigl( \SIGMA_{i} + \XI_{i} \bigr)
\end{equation}

The mixnode then forwards the updated header $\mathcal H_{i+1}=(\ALPHA_{i+1}, \mathcal{L}_{i+1}, \SIGMA_{i+1})$ to the address associated with $\ETA_{i+1}$. 
At the final hop, $\ETA_{\ell+1}=\D{}$, the destination address $d$ is recovered as $d=\operatorname{Dec}(\D{})$.


% Algorithms
% \begin{figure}
% \centering
% \begin{minipage}{0.47\linewidth}
% \begin{algorithm}[H]
% \small
% \caption{Preprocessing}% (Fig.~\ref{fig:chunk-preprocessing})}
\begin{algorithmic}[1]

\Statex \textbf{Input:}
\Statex \quad
  \begin{tabular}{@{}ll@{}}
    $(s_i)_{i=1}^{\ell}$ & shared secrets \\
    $(N_i)_{i=1}^{\ell}$ & route 
  \end{tabular}
\Statex \textbf{Output:} 
\Statex \quad
  \begin{tabular}{@{}ll@{}}
    $\PHI$ & \hspace{8.5mm} filler chunks \\
  \end{tabular}
\Statex

\State $\PHI \gets \text{empty list}$
\For{$i=1$ \textbf{to} $\ell-1$}
    \State $\PHI \gets
        \PHI \,\Vert\, H_{\mathbb{G}_1}(s_i) \,\Vert\, N_i $
    \State $\PHI \gets [s_i]\PHI$
\EndFor
\State \Return $\PHI$

\end{algorithmic}
% \label{algo:preprocessing}
% \end{algorithm}
% \end{minipage}
% \hfill
% \begin{minipage}{0.47\linewidth}
% \begin{algorithm}[H]
% \small
% \caption{Encryption}
\label{algo:encryption}
\begin{algorithmic}[1]

\Statex \textbf{Input:}
\Statex \quad
  \begin{tabular}{@{}ll@{}}
    $\D{}$                & destination point \\
    $(N_i)_{i=1}^{\ell}$   & route identifiers \\
    $(s_i)_{i=1}^{\ell}$   & shared secrets
  \end{tabular}
\Statex \textbf{Output:}
\Statex \quad
  \begin{tabular}{@{}ll@{}}
    $L_1$ & routing layer
  \end{tabular}
\Statex

\State $\beta_\ell \gets [s_\ell^{-1}]\Delta \,\Vert\, \PHI$
\State $\C{}_\ell \gets \mathsf{HMAC}(s_\ell,\beta_\ell)$
\State $L_\ell \gets (N_\ell, \C{}_\ell, \beta_\ell)$

\For{$i=\ell-1$ \textbf{down to} $1$}
    \State $\beta_i \gets [s_i^{-1}] \: L_{i+1}[1:2\ell-1]$
    \State $\C{i} \gets \mathsf{HMAC}(s_i,\beta_i)$
    \State $L_i \gets (N_i,\C{i}, \beta_i)$
\EndFor

\State \Return $L_1$

\end{algorithmic}
% \label{algo:sttp}
% \end{algorithm}
% \end{minipage}
% \end{figure}